\section{Astrobiological Cognition}\label{sec:astrobiology_cognition}
If the OBH captures a computational necessity rather than a mammalian quirk, then sleep-like alternations should emerge anywhere
complex learners evolve. Convergent cases already exist on Earth: cephalopods display REM-like ``active sleep,'' while songbirds
and aquatic mammals protect unihemispheric cycles when constant vigilance is required. These examples suggest that a dual-phase
algorithm---maintenance plus regularization---is a low-energy way to balance plasticity and stability across species \citep{hoel2021overfitted}.

\section{Would Aliens Dream?}\label{sec:aliens_dream}
Two constraints make an affirmative answer plausible. First, any intelligence that learns from finite sensory data must avoid
overfitting, implying the need for dropout- or replay-like phases whether implemented biologically or in silicon \citep{franklin2023precision}.
Second, continuous on-line inference without consolidation leads to precision drift and catastrophic rigidity, as seen in both
transformers and sleep-deprived humans \citep{chen2023attention_sink}. An extraterrestrial mind might not exhibit REM twitches, but it
would likely quarantine periods for metabolite clearance, representational renormalization, or synthetic dream generation. If we
ever detect cognition elsewhere, irregular alternations between high-fidelity sensing and internally generated, hallucinatory
simulation may be among the best biosignatures to seek.
